We introduce Distinction Graph Machines, a runtime-learning model whose state is a growing graph of anchored concepts, movable retrieval centroids, local memories, and pairwise distinctions. The guiding rule is refine when the current concepts are too coarse, and repair when a local concept is adequate but its state is wrong. The theory separates these operations. Minimal refinement is the smallest partition or sigma-algebra update that makes new evidence measurable, and contradictions inside an atom cannot be solved by any predictor measurable with respect to the current information. The value of a refinement is exact: under log loss it is conditional mutual information, and under squared loss it is conditional variance reduction. The concrete data structure realizes refinement by creating anchors and distinction edges that generate a finite information structure, while centroids and normalized edge gates provide a trainable soft read path. We give prediction and observation algorithms, monotone-repair and statistically guarded objective-descent guarantees for refine-or-repair, a data-structure complexity theorem, and a language-modeling instance in which hidden-state concepts store next-token memories and produce runtime logit corrections without inference-time backpropagation. We further define DGM-Embedding and DGM-Adapter Transformer layers with differentiable sparse reads and detached local writes, allowing slow parameters to be trained by backpropagation while the adaptive state is updated by the same forward-only rule used at inference. Projection Memory and CPM appear as local repair mechanisms after the relevant distinction structure has been fixed.
